\documentclass[a4paper,12pt]{article}

% -------------------------------------------------
% Кодировка и язык
% -------------------------------------------------
\usepackage[T2A]{fontenc}
\usepackage[utf8]{inputenc}
\usepackage[russian]{babel}

% -------------------------------------------------
% Математика
% -------------------------------------------------
\usepackage{amsmath,amssymb,amsthm,mathtools}
\usepackage{icomma}

% -------------------------------------------------
% Шрифт и типографика
% -------------------------------------------------
\usepackage{helvet}
\renewcommand{\familydefault}{\sfdefault}
\usepackage[a4paper,left=20mm,right=20mm,top=20mm,bottom=22mm]{geometry}
\usepackage{microtype}
\usepackage{parskip}
\usepackage{setspace}
\onehalfspacing

% -------------------------------------------------
% Графика
% -------------------------------------------------
\usepackage{tikz}
\usetikzlibrary{calc}
\usetikzlibrary{arrows.meta,positioning,shapes.geometric}

\usepackage{tkz-euclide}

\usepackage{pgfplots}
\pgfplotsset{compat=1.18}

% -------------------------------------------------
% Таблицы
% -------------------------------------------------
\usepackage{tabularx,booktabs,longtable}
\renewcommand{\arraystretch}{1.35}

% -------------------------------------------------
% Списки, колонтитулы, ссылки
% -------------------------------------------------
\usepackage{enumitem}
\usepackage{hyperref}
\usepackage{fancyhdr}
\usepackage{xcolor}

% -------------------------------------------------
% Палитра
% -------------------------------------------------
\definecolor{accent}{HTML}{008080}
\definecolor{darkaccent}{HTML}{004D4D}
\definecolor{textgray}{HTML}{333333}
\definecolor{softgray}{HTML}{6B7474}
\definecolor{linegray}{HTML}{C9D2D2}

\color{textgray}

\hypersetup{
    colorlinks=true,
    linkcolor=darkaccent,
    urlcolor=darkaccent
}

% -------------------------------------------------
% Колонтитулы
% -------------------------------------------------
\pagestyle{fancy}
\fancyhf{}
\setlength{\headheight}{15pt}
\fancyhead[L]{\small Статистика}
\fancyhead[R]{\small Ксения Клыкова}
\fancyfoot[C]{\small\thepage}

\renewcommand{\headrulewidth}{0.3pt}
\renewcommand{\footrulewidth}{0pt}

% -------------------------------------------------
% Заголовки
% -------------------------------------------------
\newcommand{\sectiontitle}[1]{%
    \vspace{0.4em}
    {\Large\bfseries\color{darkaccent}#1\par}
    \vspace{0.25em}
    {\color{linegray}\rule{\linewidth}{0.6pt}}
    \vspace{0.7em}
}

\newcommand{\subsectiontitle}[1]{%
    \vspace{0.5em}
    {\large\bfseries\color{accent}#1\par}
    \vspace{0.25em}
}

\newcommand{\important}[1]{%
    {\color{darkaccent}\bfseries #1}
}

% -------------------------------------------------
% TikZ: стили блок-схемы
% -------------------------------------------------
\tikzset{
    flowstart/.style={
        ellipse,
        draw=accent,
        line width=0.9pt,
        minimum width=42mm,
        minimum height=10mm,
        align=center,
        font=\small
    },
    flowaction/.style={
        rectangle,
        rounded corners=2mm,
        draw=accent,
        line width=0.9pt,
        text width=96mm,
        minimum height=11mm,
        align=center,
        inner sep=3mm,
        font=\small
    },
    flowactionsmall/.style={
        rectangle,
        rounded corners=2mm,
        draw=accent,
        line width=0.9pt,
        text width=55mm,
        minimum height=13mm,
        align=center,
        inner sep=3mm,
        font=\small
    },
    flowdecision/.style={
        diamond,
        aspect=2.5,
        draw=accent,
        line width=0.9pt,
        text width=44mm,
        align=center,
        inner sep=1.5mm,
        font=\small
    },
    flowarrow/.style={
        -{Stealth[length=2.5mm,width=1.8mm]},
        line width=0.9pt,
        draw=darkaccent
    }
}

% -------------------------------------------------
% Документ
% -------------------------------------------------
\begin{document}

% =================================================
% ТИТУЛЬНЫЙ ЛИСТ
% =================================================
\begin{titlepage}
\thispagestyle{empty}

\begin{center}

\vspace*{24mm}

{\large\color{softgray} ЕГЭ по профильной математике\par}

\vspace{18mm}

{\Huge\bfseries\color{darkaccent} Чек-лист\par}

\vspace{5mm}

{\LARGE\bfseries
Статистические данные:\\[2mm]
частоты, выборка и средние значения
\par}

\vspace{16mm}

{\large
Ксения Клыкова
\par}

\vspace{8mm}

{\normalsize
Вариационный ряд \(\cdot\)
абсолютная и относительная частоты \(\cdot\)
объём выборки
\par}

\vspace{3mm}

{\normalsize
таблицы частот \(\cdot\)
среднее арифметическое \(\cdot\)
взвешенное среднее
\par}

\vfill

{\large \today\par}

\vspace{10mm}

{\normalsize
\textbf{Лёвин Артём Александрович}\\[2mm]
эксклюзивно для Ксении Клыковой
\par}

\vspace{14mm}

\begin{tikzpicture}[remember picture,overlay]
    \draw[
        accent,
        line width=1.1pt
    ]
    ($(current page.south west)+(24mm,19mm)$)
    .. controls
    ($(current page.south west)+(62mm,31mm)$)
    and
    ($(current page.south east)+(-65mm,7mm)$)
    ..
    ($(current page.south east)+(-24mm,19mm)$);
\end{tikzpicture}

\end{center}
\end{titlepage}

% =================================================
% ОСНОВНОЙ ЧЕК-ЛИСТ
% =================================================

\sectiontitle{1. Что нужно знать}

\subsectiontitle{Вариационный ряд}

\important{Вариационный ряд} --- числовые данные, расположенные в порядке
неубывания.

Например, для исходного набора
\[
5,\;2,\;4,\;2,\;3,\;4,\;2
\]
вариационный ряд имеет вид
\[
2,\;2,\;2,\;3,\;4,\;4,\;5.
\]

Вариационный ряд помогает быстро увидеть повторяющиеся значения, минимальное
и максимальное значения, а также подготовить данные к составлению таблицы
частот.

\subsectiontitle{Объём выборки}

\important{Объём выборки} \(n\) --- общее количество наблюдений.

Для набора
\[
2,\;2,\;2,\;3,\;4,\;4,\;5
\]
имеем
\[
n=7.
\]

Важно учитывать \emph{каждое} наблюдение, включая повторяющиеся значения.

\subsectiontitle{Абсолютная частота}

\important{Абсолютная частота} \(m_i\) показывает, сколько раз конкретное
значение \(x_i\) встретилось в выборке.

Для приведённого выше ряда:
\[
m(2)=3,\qquad
m(3)=1,\qquad
m(4)=2,\qquad
m(5)=1.
\]

Сумма всех абсолютных частот всегда совпадает с объёмом выборки:
\[
\boxed{\sum_{i=1}^{k}m_i=n}.
\]

Здесь \(k\) --- количество различных значений признака.

\subsectiontitle{Относительная частота}

\important{Относительная частота} \(p_i\) показывает, какую долю всей выборки
составляют наблюдения со значением \(x_i\).

Формула:
\[
\boxed{p_i=\frac{m_i}{n}}.
\]

Для значения \(2\):
\[
p(2)=\frac{3}{7}.
\]

Обратная формула:
\[
\boxed{m_i=p_i n}.
\]

Сумма всех относительных частот равна единице:
\[
\boxed{\sum_{i=1}^{k}p_i=1}.
\]

Если относительные частоты записаны в процентах, их сумма равна \(100\%\).

\subsectiontitle{Почему относительная частота полезна}

Само число повторений ещё требуется соотнести с объёмом выборки.

Например, значение встретилось \(3\) раза.

Если
\[
n=7,
\]
то его доля составляет
\[
\frac{3}{7}\approx 0{,}429.
\]

Если
\[
n=100,
\]
то его доля составляет всего
\[
\frac{3}{100}=0{,}03.
\]

Поэтому относительная частота позволяет сравнивать представленность значения
в выборках разного объёма.

% =================================================

\sectiontitle{2. Таблица частот}

Рассмотрим вариационный ряд
\[
1,\;2,\;2,\;3,\;3,\;3,\;4,\;4.
\]

Объём выборки:
\[
n=8.
\]

\begin{table}[ht]
\centering
\caption{Таблица абсолютных и относительных частот}
\begin{tabularx}{0.88\linewidth}{>{\centering\arraybackslash}X
                                >{\centering\arraybackslash}X
                                >{\centering\arraybackslash}X}
\toprule
Значение \(x_i\) &
Абсолютная частота \(m_i\) &
Относительная частота \(p_i\) \\
\midrule
\(1\) & \(1\) & \(\frac18\) \\
\(2\) & \(2\) & \(\frac28=\frac14\) \\
\(3\) & \(3\) & \(\frac38\) \\
\(4\) & \(2\) & \(\frac28=\frac14\) \\
\midrule
Итого & \(8\) & \(1\) \\
\bottomrule
\end{tabularx}
\end{table}

При работе с таблицей полезно выполнять две контрольные проверки:
\[
\sum m_i=n,
\qquad
\sum p_i=1.
\]

Если одна из проверок нарушается, в вычислениях или при переносе данных
появилась ошибка.

\subsectiontitle{Пример: неизвестная частота}

Значения \(0,\;1,\;2,\;3\) имеют абсолютные частоты
\[
4,\;8,\;a,\;3.
\]
Всего выполнено \(20\) наблюдений.

Используем
\[
\sum m_i=n:
\]
\[
4+8+a+3=20.
\]

Отсюда
\[
a=20-15=5.
\]

\important{Приём:} если известен объём выборки и все частоты, кроме одной,
неизвестную частоту удобно находить через сумму абсолютных частот.

% =================================================

\sectiontitle{3. Интервальные данные}

Значением признака может быть целый интервал.

Например:

\begin{table}[ht]
\centering
\caption{Распределение времени выполнения задания}
\begin{tabularx}{0.78\linewidth}{>{\centering\arraybackslash}X
                                >{\centering\arraybackslash}X}
\toprule
Время, мин & Число учеников \\
\midrule
от \(0\) до \(10\) & \(6\) \\
от \(10\) до \(20\) & \(14\) \\
от \(20\) до \(30\) & \(10\) \\
\midrule
Итого & \(30\) \\
\bottomrule
\end{tabularx}
\end{table}

Требуется определить долю учеников, справившихся менее чем за \(20\) минут.

Под условие подходят две первые группы:
\[
6+14=20.
\]

Поэтому искомая доля равна
\[
\frac{20}{30}=\frac23.
\]

\important{Типичная особенность:} требуемое условие может охватывать несколько
строк или столбцов таблицы. Тогда сначала складывают соответствующие частоты.

% =================================================

\sectiontitle{4. Среднее арифметическое}

Для чисел
\[
x_1,x_2,\ldots,x_n
\]
среднее арифметическое обозначают \(\overline{x}\) и вычисляют по формуле
\[
\boxed{
\overline{x}
=
\frac{x_1+x_2+\ldots+x_n}{n}
}.
\]

Компактная запись:
\[
\boxed{
\overline{x}
=
\frac{1}{n}\sum_{i=1}^{n}x_i
}.
\]

\subsectiontitle{Пример}

Пять результатов равны
\[
10,\;12,\;14,\;16,\;18.
\]

Тогда
\[
\overline{x}
=
\frac{10+12+14+16+18}{5}
=
\frac{70}{5}
=
14.
\]

Для любого набора чисел среднее значение располагается между наименьшим и
наибольшим наблюдаемыми значениями:
\[
x_{\min}\le \overline{x}\le x_{\max}.
\]

Это полезная проверка ответа.

% =================================================

\sectiontitle{5. Среднее по таблице частот}

Если одно и то же значение встречается много раз, каждое повторение отдельно
выписывать неудобно.

Пусть значения
\[
x_1,x_2,\ldots,x_k
\]
имеют абсолютные частоты
\[
m_1,m_2,\ldots,m_k.
\]

Тогда
\[
\boxed{
\overline{x}
=
\frac{x_1m_1+x_2m_2+\ldots+x_km_k}
     {m_1+m_2+\ldots+m_k}
}.
\]

Поскольку
\[
m_1+m_2+\ldots+m_k=n,
\]
формулу можно записать компактно:
\[
\boxed{
\overline{x}
=
\frac{\displaystyle\sum_{i=1}^{k}x_i m_i}{n}
}.
\]

Такое вычисление часто называют \important{взвешенным средним}.

\subsectiontitle{Пример: средняя цена}

За день продали:

\begin{itemize}[leftmargin=8mm,itemsep=2pt]
    \item \(10\) тетрадей по \(30\) рублей;
    \item \(20\) тетрадей по \(35\) рублей.
\end{itemize}

Всего продано
\[
10+20=30
\]
тетрадей.

Средняя цена одной проданной тетради:
\[
\overline{x}
=
\frac{30\cdot10+35\cdot20}{30}
=
\frac{1000}{30}
=
\frac{100}{3}
=
33{,}(3)\text{ руб.}
\]

Простое выражение
\[
\frac{30+35}{2}
\]
здесь не учитывает различное количество тетрадей каждой цены.

% =================================================

\sectiontitle{6. Задачи на изменение среднего}

Если известно среднее значение и число наблюдений, сначала удобно восстановить
их сумму:
\[
\boxed{
S=n\overline{x}
}.
\]

\subsectiontitle{Пример: удалили одно число}

Среднее арифметическое \(8\) чисел равно \(12\).

Их общая сумма:
\[
S_1=8\cdot12=96.
\]

После удаления одного числа среднее арифметическое оставшихся \(7\) чисел стало
равно \(10\).

Их сумма:
\[
S_2=7\cdot10=70.
\]

Удалённое число равно разности сумм:
\[
96-70=26.
\]

Ответ:
\[
\boxed{26}.
\]

\important{Ключевая идея:}
\[
\text{сумма данных}
=
\text{число наблюдений}\cdot
\text{среднее}.
\]

После этого изменение выборки удобно анализировать через изменение её суммы.



% =================================================
% ЧЕК-ЛИСТ ОШИБОК
% =================================================

\sectiontitle{7. Что проверять перед ответом}

\begin{enumerate}[leftmargin=9mm,itemsep=4pt]
    \item
    \textbf{Объём выборки.}
    Все наблюдения, включая повторяющиеся, должны быть учтены.

    \item
    \textbf{Абсолютные частоты.}
    Выполните проверку
    \[
    \sum m_i=n.
    \]

    \item
    \textbf{Относительные частоты.}
    Выполните проверку
    \[
    \sum p_i=1.
    \]

    \item
    \textbf{Интервалы.}
    Проверьте, какие интервалы полностью удовлетворяют условию задачи.

    \item
    \textbf{Среднее арифметическое.}
    В знаменателе должно стоять количество наблюдений.

    \item
    \textbf{Повторяющиеся значения.}
    При вычислении среднего по таблице учитывайте частоты:
    \[
    x_i m_i.
    \]

    \item
    \textbf{Изменение среднего.}
    Восстановите сумму по формуле
    \[
    S=n\overline{x}.
    \]

    \item
    \textbf{Оценка результата.}
    Для обычного среднего должно выполняться
    \[
    x_{\min}\le\overline{x}\le x_{\max}.
    \]
\end{enumerate}

\subsectiontitle{Мини-шпаргалка}

\[
\boxed{
p_i=\frac{m_i}{n}
}
\qquad
\boxed{
m_i=p_i n
}
\]

\[
\boxed{
\sum_{i=1}^{k}m_i=n
}
\qquad
\boxed{
\sum_{i=1}^{k}p_i=1
}
\]

\[
\boxed{
\overline{x}
=
\frac{\sum x_i}{n}
}
\qquad
\boxed{
\overline{x}
=
\frac{\sum x_i m_i}{n}
}
\]

\[
\boxed{
S=n\overline{x}
}
\]

% =================================================
% ТРЕНИРОВОЧНЫЙ БЛОК — ОТДЕЛЬНАЯ СТРАНИЦА
% =================================================

\clearpage

\sectiontitle{Тренировочный блок: 5 заданий}

\textbf{Задание 1. Вариационный ряд и частоты}

Дан набор:
\[
4,\;2,\;4,\;1,\;2,\;4,\;3,\;2,\;5,\;2.
\]

\begin{enumerate}[label=\alph*),leftmargin=9mm,itemsep=1pt]
    \item составьте вариационный ряд;
    \item найдите объём выборки;
    \item найдите абсолютную частоту значения \(2\);
    \item найдите относительную частоту значения \(2\).
\end{enumerate}

\vspace{5mm}

\textbf{Задание 2. Таблица частот}

Время пути учеников до школы распределилось следующим образом:

\begin{center}
\begin{tabular}{ccccc}
\toprule
Время, мин & \(10\) & \(15\) & \(20\) & \(25\) \\
\midrule
Число учеников & \(3\) & \(5\) & \(4\) & \(8\) \\
\bottomrule
\end{tabular}
\end{center}

Найдите:

\begin{enumerate}[label=\alph*),leftmargin=9mm,itemsep=1pt]
    \item объём выборки;
    \item относительную частоту каждого значения;
    \item сумму всех относительных частот.
\end{enumerate}

\vspace{5mm}

\textbf{Задание 3. Неизвестная абсолютная частота}

Значения
\[
0,\;1,\;2,\;3
\]
имеют абсолютные частоты
\[
4,\;8,\;a,\;3.
\]

Всего выполнено \(20\) наблюдений.

Найдите:

\begin{enumerate}[label=\alph*),leftmargin=9mm,itemsep=1pt]
    \item \(a\);
    \item относительную частоту значения \(2\).
\end{enumerate}

\vspace{5mm}

\textbf{Задание 4. Интервальные данные}

Время выполнения задания распределилось так:

\begin{center}
\begin{tabular}{cccc}
\toprule
Время, мин & \(0\)--\(10\) & \(10\)--\(20\) & \(20\)--\(30\) \\
\midrule
Число учеников & \(7\) & \(11\) & \(12\) \\
\bottomrule
\end{tabular}
\end{center}

Какую долю всех учеников составляют ученики, выполнившие работу менее чем за
\(20\) минут?

\vspace{5mm}

\textbf{Задание 5. Взвешенное среднее}

В магазине продали \(12\) одинаковых по типу тетрадей по \(48\) рублей и
\(18\) тетрадей по \(56\) рублей.

Найдите среднюю цену одной проданной тетради.

% =================================================
% ОТВЕТЫ — ОТДЕЛЬНАЯ СТРАНИЦА
% =================================================

\clearpage

\sectiontitle{Ответы к тренировочному блоку}

\begin{table}[ht]
\centering
\caption{Ответы}
\begin{tabularx}{\linewidth}{>{\centering\arraybackslash}p{18mm} X}
\toprule
№ & Ответ \\
\midrule

1 &
\(
1,\;2,\;2,\;2,\;2,\;3,\;4,\;4,\;4,\;5;
\quad
n=10;
\quad
m(2)=4;
\quad
p(2)=\frac{4}{10}=\frac25=0{,}4.
\)
\\

\addlinespace[2mm]

2 &
\(
n=20.
\)
Относительные частоты:
\[
\frac{3}{20}=0{,}15,\qquad
\frac{5}{20}=0{,}25,\qquad
\frac{4}{20}=0{,}20,\qquad
\frac{8}{20}=0{,}40.
\]
Их сумма равна \(1\).
\\

\addlinespace[2mm]

3 &
\[
4+8+a+3=20,
\qquad
a=5.
\]
\[
p(2)=\frac{5}{20}=\frac14=0{,}25.
\]
\\

\addlinespace[2mm]

4 &
\[
\frac{7+11}{7+11+12}
=
\frac{18}{30}
=
\frac35.
\]
\\

\addlinespace[2mm]

5 &
\[
\overline{x}
=
\frac{48\cdot12+56\cdot18}{12+18}
=
\frac{1584}{30}
=
52{,}8.
\]
Средняя цена --- \(52{,}8\) руб.
\\

\bottomrule
\end{tabularx}
\end{table}

\vfill

\begin{center}
{\color{softgray}\small
Главная цель блока --- уверенно переводить текст статистической задачи
в значения, частоты, объём выборки и подходящую формулу.
}
\end{center}

\end{document}